\documentclass[11pt,a4paper]{article}

\usepackage{amsmath,amssymb,amsthm}
\usepackage[margin=2.5cm]{geometry}
\usepackage{mathtools}
\usepackage{hyperref}
\usepackage{microtype}
\usepackage{graphicx}
\usepackage{booktabs}

\newcommand{\norm}[1]{\left\|#1\right\|}
\newcommand{\abs}[1]{\left|#1\right|}
\newcommand{\C}{\mathbb{C}}
\newcommand{\spec}{\sigma}
\newcommand{\diag}{\operatorname{diag}}
\DeclareMathOperator{\ad}{ad}

\theoremstyle{plain}
\newtheorem{theorem}{Theorem}[section]
\newtheorem{lemma}[theorem]{Lemma}
\newtheorem{corollary}[theorem]{Corollary}
\newtheorem{proposition}[theorem]{Proposition}
\theoremstyle{definition}
\newtheorem{remark}[theorem]{Remark}
\newtheorem{definition}[theorem]{Definition}

\title{\textbf{A certified resolvent bound via verified block diagonalization}\\[0.3em]
\large A block-orthonormal Schur frame, a Miyajima-style proof, and the $\sigma_{\min}$ surface}
\author{}
\date{\today}

\begin{document}
\maketitle

\begin{abstract}
\noindent
Given a dense matrix $A\in\C^{n\times n}$, we produce a rigorous, cheaply swept upper
bound on the resolvent norm $\norm{(A-zI)^{-1}}_2$ over the complex plane ---
equivalently, a certified lower bound on $\sigma_{\min}(A-zI)$. The bound is read off a
\emph{verified block-diagonal decomposition} $A\approx W\,\operatorname{blockdiag}(P_j)\,W^{-1}$
(Miyajima's verified block diagonalization): a cheap numerical phase builds the frame $W$
and the blocks from a Schur factorization followed by a guarded Newton--Rump separation,
and a single rigorous phase certifies the result through Miyajima's two-residual argument,
which never forms an exact inverse and folds all conditioning into one residual norm. The
frame is \emph{block-orthonormal} --- orthonormal within each block but not between --- so
its conditioning $\kappa(W)$ is governed by benign inter-block subspace angles, while the
malignant intra-block non-normality (the defective or near-defective part) stays inside
$P_j$, where a verified $\sigma_{\min}(P_j-zI)$ absorbs it. The resolvent bound is then
$\min_j\sigma_{\min}(P_j-zI)$ minus the certified inter-block coupling, divided by
$\kappa(W)$; a spectral enclosure --- discs whose union contains $\spec(A)$, with each
cluster assigned exactly as many eigenvalues as it holds --- falls out as the locus where
the bound degenerates. We give the geometric picture (the off-diagonal of $W^{-1}AW$ is the
``leakage'' of $A$ along the frame, and the Newton step is the reduced resolvent that
removes it), the algorithm, the Miyajima proof, and experiments, including the
Gauss--Kuzmin--Wirsing transfer operator in extended precision. Finally we extend the
method to operators available \emph{only through matrix--vector products}: the near-target
eigenvalues are captured in a small dense block, and the complementary far block is
certified by a projected Lanczos process whose loss of orthogonality is bounded
\emph{a posteriori} from scalar coefficients via Paige's recurrence --- with no stored
basis. Run to a full basis, the single Neumann residual $\norm{R_2}_2<1$ certifies at once
that the Lanczos frame is invertible and that it is complete, so no eigenvalue escapes the
enclosure; far from the spectrum, where the method is used, it is sound and fires.
\end{abstract}

\section{Setup}\label{sec:setup}

Let $A\in\C^{n\times n}$. A numerical eigensolver returns an approximate
eigenvector matrix $W\in\C^{n\times n}$ and approximate eigenvalues $d_1,\dots,d_n$,
with $D=\diag(d_1,\dots,d_n)$, so that $AW\approx WD$. We do \emph{not} assume $W$
is exactly invertible; nonsingularity will be \emph{verified}
(Section~\ref{sec:rigor}). Where $W$ is invertible, the \emph{collapsed matrix}
\begin{equation}\label{eq:collapse}
  M \coloneqq W^{-1}AW
\end{equation}
is similar to $A$, so $\spec(M)=\spec(A)$, and is \emph{almost diagonal}:
\[
  M = D + E,\qquad E \coloneqq M-D \ \text{(small off-diagonal residue)}.
\]
Everything that follows reads the spectrum off $M$; the off-diagonal $E$ measures the
failure of $W$ to diagonalize $A$, and the whole method is about (i) shrinking $E$
cheaply (Section~\ref{sec:geom}--\ref{sec:algo}) and (ii) bounding the eigenvalues of
$M$ rigorously from a computed proxy of it (Section~\ref{sec:rigor}).

\section{Geometric intuition: leakage and the reduced resolvent}\label{sec:geom}

Write $w_1,\dots,w_n$ for the columns of $W$ (the approximate eigenvectors). From
$AW=WM$, the $j$-th column reads
\begin{equation}\label{eq:leakage}
  A w_j \;=\; \sum_{i} M_{ij}\, w_i \;=\; \underbrace{d_j\, w_j}_{\text{eigen part}}
            \;+\; \underbrace{\sum_{i\ne j} M_{ij}\, w_i}_{\text{leakage}\,=\,r_j}.
\end{equation}
Thus $M_{ij}$ ($i\ne j$) is exactly the $w_i$-component of $Aw_j$: the amount by which
$Aw_j$ ``leaks'' out of the direction $w_j$ into the other eigendirections. For an
exact eigenvector the leakage vanishes and $M$ is diagonal; the residual
$r_j = Aw_j-d_jw_j = \sum_{i\ne j}M_{ij}w_i$ quantifies how far $w_j$ is from being an
eigenvector.

To improve $w_j$ we tilt it toward the other approximate eigenvectors,
$w_j' = w_j + \sum_{i\ne j} X_{ij}\,w_i$, i.e. $W' = W(I+X)$ with $X$ off-diagonal,
and ask the tilted vector to be (closer to) an eigenvector. To first order the
$w_i$-component of $A w_j'$ is $M_{ij} + d_i X_{ij}$, and we want it to equal
$d_j X_{ij}$ (so $Aw_j'\approx d_j w_j'$). Solving,
\begin{equation}\label{eq:correction}
  M_{ij} + d_i X_{ij} = d_j X_{ij}
  \quad\Longrightarrow\quad
  \boxed{\,X_{ij} = -\frac{M_{ij}}{d_i-d_j}\,}\qquad(i\ne j).
\end{equation}

\paragraph{This is the reduced resolvent.} Equation~\eqref{eq:correction} is the
action, in the eigenbasis, of the resolvent of $A$ at $d_j$ restricted to the other
eigendirections: along $w_i$ the operator $(d_jI-A)$ scales by $d_j-d_i$, so its
inverse divides by $d_j-d_i$. Writing $S_j = (A-d_jI)^{-1}\big|_{\operatorname{span}(w_{i\ne j})}$
for that \emph{reduced resolvent}, the update is
\[
  w_j' \;=\; w_j \;-\; S_j\, r_j,
\]
the classical eigenvector Newton (equivalently first-order Rayleigh--Schr\"odinger)
step: \emph{forward} the eigenvector through $A$ to expose the leakage $r_j$, then
\emph{pull it back} through the reduced resolvent --- divide each component by its
spectral gap $d_i-d_j$.

\paragraph{Sylvester form.} Collecting \eqref{eq:correction} over all $i,j$, killing
the off-diagonal of $M$ to first order is the Sylvester equation
\begin{equation}\label{eq:sylvester}
  [D,X] = DX-XD = -E,\qquad\text{i.e.}\qquad (d_i-d_j)X_{ij} = -E_{ij}.
\end{equation}
The commutator operator $\ad_D = [D,\cdot]$ is diagonal in the entrywise basis with
eigenvalues $d_i-d_j$, so its inverse is the Hadamard division by the gaps:
$X = -\ad_D^{-1}(E)$, $X_{ij}=-E_{ij}/(d_i-d_j)$. The matrix $1/(d_i-d_j)$ is the
``reciprocal-distance preconditioner'' of the Newton iteration.

\paragraph{Two consequences fall out of the geometry.}
\begin{itemize}
  \item \emph{Quadratic convergence.} The update keeps only the first-order term; the
    new collapsed matrix $M'=(I+X)^{-1}M(I+X)$ has off-diagonal $O(\norm{E}^2)$, so a
    sweep of $w_j' = w_j - S_jr_j$ squares the residual. Iterating
    \eqref{eq:correction} (with $M$ recomputed) is the refinement of
    Section~\ref{sec:algo}.
  \item \emph{Coincident eigenvalues are forbidden, and that is correct.} When
    $d_i=d_j$ the gap closes and $S_j$ is singular: there is no direction to tilt $w_j$
    ``toward'' inside a degenerate eigenspace, where eigenvectors are not uniquely
    defined. Numerically the step diverges once $\abs{M_{ij}}\gtrsim\abs{d_i-d_j}$ (it
    no longer contracts). The fix is to refuse the step there --- to \emph{cluster}
    $i,j$ and leave them coupled. This is the only place clustering enters, and it is
    dictated by the geometry, not by a modulus heuristic.
\end{itemize}

\section{The algorithm}\label{sec:algo}

The method separates a cheap \emph{numerical} phase --- which produces a candidate
\emph{block-diagonal decomposition} of $A$ --- from a single \emph{rigorous} phase (the
certification of Section~\ref{sec:rigor}). Correctness rests only on the latter, so all
refinement and clustering decisions may be taken in plain floating point.

\paragraph{What is produced.} We do not attempt the full diagonalization: for a defective
or near-defective $A$ the eigenbasis is ill-conditioned (or does not exist), and forcing
it inflates $\kappa(W)$ by $u^{-1/k}$ on a $k$-fold cluster. The numerical phase returns
instead a \emph{block}-diagonal decomposition
\[
  A \;\approx\; W\,\Lambda\,W^{-1},\qquad
  \Lambda=\operatorname{blockdiag}(P_1,\dots,P_q),
\]
in which well-separated eigenvalues become $1\times1$ blocks and eigenvalues too close to
separate at the working precision are kept together as one block $P_j$. This object
always exists, and it is built so that $W$ is \emph{block-orthonormal}: its columns are
orthonormal \emph{within} each block, so $\kappa(W)$ is governed by the \emph{inter-block}
subspace angles ($\approx 1/\sin\theta_{\mathrm{inter}}$, benign and separation-controlled)
and never by intra-block eigenvector near-parallelism. That malignant non-normality stays
\emph{inside} the block $P_j$, where the verified $\sigma_{\min}(P_j-zI)$ absorbs it
(Section~\ref{sec:resolvent}), rather than poisoning the global prefactor. The full
diagonalization is the special case in which every block is $1\times1$.

\paragraph{The approximate inverse.} Following Miyajima, we never form $W^{-1}$.
Instead we take \emph{any} computed $Y\approx W^{-1}$ (e.g.\ a floating-point inverse
of $W$) and measure its quality by the single residual
\begin{equation}\label{eq:R2}
  R_2 \coloneqq YW - I .
\end{equation}
The certificate is valid for any $Y$ with $\norm{R_2}<1$ (Lemma~\ref{lem:neumann}):
$Y$ is free, and its inaccuracy enters only through $\norm{R_2}$.

\paragraph{The pipeline.} Four steps; only the first is $O(n^3)$ beyond the per-step
matrix products.
\begin{enumerate}
  \item \textbf{Schur.} $A=QTQ^{H}$ (\texttt{GenericSchur} in extended precision), $Q$
    \emph{orthonormal}, $T$ upper triangular with the eigenvalues on its diagonal. The
    orthonormal Schur vectors are the well-conditioned starting frame, and the collapsed
    matrix $M=Q^{H}AQ=T$ is already almost block-triangular --- so the geometry of
    Section~\ref{sec:geom} applies with $W=Q$.
  \item \textbf{Newton--Rump separation.} Diagonalize \emph{between} clusters and stop
    \emph{within} them: apply the reduced-resolvent step \eqref{eq:correction} to the
    cross-cluster entries only, $W\leftarrow W(I+X)$, $X_{ij}=-M_{ij}/(d_i-d_j)$ for $i,j$
    in different clusters. The clustering is fixed by a \emph{merge-only,
    self-correcting} rule, after Rump's \texttt{verifyeigall}: start from the finest
    grouping (singletons separated by a distance threshold $\tau\sim\sqrt u\,\norm A$) and
    \emph{merge} any column that fails to contract --- the floating-point surrogate
    ($\norm{X_{:,j}}_2\ge 1$, or the precision floor $\abs{d_i-d_j}<u^{1/k}\norm A$ for a
    tentative $k$-block) for Rump's rigorous verification failing --- iterating until no
    column fails. Clusters only \emph{grow}; this is the clustering of
    Section~\ref{sec:geom} (refuse the step where it does not contract), made
    self-correcting. A genuinely near-defective cluster --- e.g.\ a geometrically decaying
    tail --- coalesces into one block rather than being split into ill-posed pieces. The
    per-pair test would miss the \emph{cumulative} coupling of such a chain; the
    column-norm test sees it.
  \item \textbf{Block-diagonal approximation.} On the separated $M=W^{-1}AW$, zero the
    (now small) off-block entries to read off the candidate
    $\Lambda=\operatorname{blockdiag}(P_1,\dots,P_q)$, $P_j=M_{\mathcal C_j\mathcal C_j}$.
    The discarded off-block part is the inter-block coupling, certified and carried as
    $\norm C$ in Sections~\ref{sec:rigor}--\ref{sec:resolvent}.
  \item \textbf{Orthogonalization inside the blocks.} Re-orthonormalize the columns of
    $W$ \emph{within} each block ($W_{:,\mathcal C_j}\leftarrow$ the $Q$-factor of its thin
    QR). Because the blocks are decoupled (Step 3) this touches nothing else, and it yields
    the block-orthonormal frame above. We do \emph{not} orthonormalize \emph{between}
    blocks: the invariant subspaces of a non-normal $A$ are not mutually orthogonal, so
    doing so would rotate the (exactly block-diagonal) $\Lambda$ back to
    block-\emph{triangular} --- reintroducing the very coupling Step~2 removed. The inter-block
    conditioning is instead carried honestly through $\kappa(W)$ in the certificate.
\end{enumerate}
The geometry of Section~\ref{sec:geom} supplies both the step (the reduced resolvent
\eqref{eq:correction}, quadratically convergent between separated blocks) and the
clustering (refuse it where the gap closes); the VBD pipeline only chooses the orthonormal
Schur frame and confines the residual non-normality to the diagonal blocks.

\begin{remark}[Block refinement and regional queries]\label{rem:block}
Step~2 applies \eqref{eq:correction} entrywise, but the same correction has a block
form. For any partition of the indices into blocks, the coupling between two blocks
$\mathcal B,\mathcal B'$ is removed to first order by $W\leftarrow W(I+X)$ with $X$
solving the \emph{Sylvester equation}
\[
  M_{\mathcal B\mathcal B}\,X - X\,M_{\mathcal B'\mathcal B'} = -\,M_{\mathcal B\mathcal B'},
\]
whose entrywise solution is again $X_{ij}=-M_{ij}/(d_i-d_j)$, $i\in\mathcal B$,
$j\in\mathcal B'$ --- \eqref{eq:correction} on cross-block pairs, well-posed when the
two blocks' spectra are separated. The working clustering is exactly such a partition
(neutralize within, decouple across), so this is not a new method but a regrouping of
the work, worth testing for efficiency: at each refinement step, choose the block
partition that best separates the current spectrum (by the contraction criterion
$\abs{M_{ij}}\ge\theta\abs{d_i-d_j}$, i.e.\ a cut of the coupling graph) and apply one
block Sylvester solve per pair of blocks, rather than a single entrywise sweep. A
region of interest is then just one distinguished block: solving across its boundary
decouples it, and only its $W$-columns need be refined --- no global refinement. The
certification (Section~\ref{sec:rigor}) is unchanged in every case: the enclosure of
Theorem~\ref{thm:enclosure} read off the relevant rows, with the count from
Corollary~\ref{cor:clusters}.
\end{remark}

\begin{remark}[Which correction? A numerical experiment]\label{rem:exp}
We compared the two corrections on matrices carrying a near-defective working cluster
(a perturbed $4\times4$ Jordan block, $\kappa_2(W)\sim10^4$, embedded with
well-separated eigenvalues), refining from $40$ randomly perturbed starting bases.
Three facts emerged. (i) Both keep \emph{every} intermediate certificate rigorous, and
they \emph{coincide} away from coupled clusters (all working blocks singletons) and at
the $\kappa\,u$ fixed point --- they differ only in the transient on the coupled
cluster. (ii) There \emph{neither dominates}: each converged on $33$--$34$ of the $40$
bases, with \emph{partly disjoint} failure sets. The entrywise sweep stalls at a
non-diagonal fixed point precisely when its cross-cluster step --- which uses only the
scalar gaps and so ignores the internal coupling of $M_{\mathcal B\mathcal B}$ ---
cannot drive the cross-coupling below the within-cluster level; the block step, exact
to first order on that coupling, escapes those stalls but has its own. (iii) A stall is
a \emph{robust attractor} of the single-variant iteration: \emph{alternating} the two
steps does not help ($33/40$, a third map with its own basins), and \emph{re-perturbing
$W$ and continuing with the same step} is reabsorbed unless the kick is tuned (a small
kick is pulled back; a $\sim\!10\%$ kick recovers $36/40$; a large kick lands in
another stall). What reliably escapes is \emph{changing the map}: run one correction to
a stall and retry with the other. Because an entrywise stall is generically not a block
stall and vice versa, this fall-back recovers the disjoint failures
($36/40$), parameter-free. The few bases that stall for \emph{both} ($4/40$ here) are
first-order--obstructed --- the cluster is too coupled for any single Newton step ---
and call for the extended-precision local re-solve (\S To~do), not a different
first-order correction. \emph{Practical rule:} default to the cheap entrywise sweep and,
on a detected stall, switch to the block-Sylvester step (or, with one variant only,
apply a moderate random kick); a per-step \emph{combination} of the two buys nothing.
\end{remark}

\paragraph{Certification.} On the final $W$, with the block-diagonal candidate
$\Lambda=\operatorname{blockdiag}(P_1,\dots,P_q)$, compute the two residuals
\begin{equation}\label{eq:Mtilde}
  R_2 = YW-I,\qquad R_1 \coloneqq Y\,(AW-W\Lambda),\qquad \widetilde M\coloneqq YAW ,
\end{equation}
verify $\norm{R_2}_2<1$, and return the enclosure of Theorem~\ref{thm:enclosure} together
with the clustering of Corollary~\ref{cor:clusters}. The block residual $R_1$ measures the
deviation of the collapsed matrix from the block-diagonal candidate; by \eqref{eq:blockdev}
the single slack $\beta_\Lambda=\norm{R_1}_2/(1-\norm{R_2}_2)$ bounds it entirely (the
inter-block leakage included --- there is \emph{no} separate coupling term), and $R_2$ both
verifies $W$ nonsingular and supplies the prefactor $\kappa_2(W)$ that lifts the bound to
$A$ (Section~\ref{sec:resolvent}). The enclosure is then $\sigma_{\min}(P_j-zI)\le\beta_\Lambda$
per block --- a recentred disc for a singleton, the block region otherwise. This is the only
step that must be performed in outward-rounded interval/ball arithmetic; in practice the
products $YW$, $YAW$ and $Y(AW-W\Lambda)$ are formed as ball matrices so that every
floating-point rounding error is rigorously absorbed.

\section{Rigorous enclosure: the Miyajima argument}\label{sec:rigor}

Throughout, $\norm{\cdot}_\infty$ is the matrix $\infty$-norm (maximal absolute row
sum) and $\spec(\cdot)$ the spectrum. The computed proxy of the collapsed matrix is
$\widetilde M = YAW$ from \eqref{eq:Mtilde}; the true collapsed matrix is $M=W^{-1}AW$
from \eqref{eq:collapse}.

\begin{lemma}[Neumann / approximate inverse]\label{lem:neumann}
Let $R_2=YW-I$ and suppose $\norm{R_2}<1$ in some submultiplicative norm. Then $W$ and
$Y$ are nonsingular, $YW=I+R_2$ is invertible with
\[
  \norm{(I+R_2)^{-1}}\le\frac{1}{1-\norm{R_2}},\qquad
  \norm{(I+R_2)^{-1}-I}\le\frac{\norm{R_2}}{1-\norm{R_2}},
\]
and $\norm{W^{-1}}\le \norm{Y}/(1-\norm{R_2})$ --- all without forming any inverse.
\end{lemma}

\begin{proof}
The Neumann series $\sum_{k\ge0}(-R_2)^k$ converges (as $\norm{R_2}<1$) to
$(I+R_2)^{-1}$, giving the first bound; subtracting $I$ leaves
$\sum_{k\ge1}(-R_2)^k$, whose norm is $\le\norm{R_2}/(1-\norm{R_2})$. Since
$I+R_2=YW$ is invertible, $W$ and $Y$ are nonsingular, and
$W^{-1}=(YW)^{-1}Y=(I+R_2)^{-1}Y$, whence $\norm{W^{-1}}\le\norm{Y}/(1-\norm{R_2})$.
\end{proof}

\begin{lemma}[The collapse is a certified perturbation of $\widetilde M$]\label{lem:deviation}
If $\norm{R_2}_\infty<1$ then $W$ is nonsingular, $\spec(A)=\spec(M)$, and
\[
  M=(I+R_2)^{-1}\widetilde M,\qquad
  \norm{M-\widetilde M}_\infty \;\le\; \beta \coloneqq
  \frac{\norm{R_2}_\infty}{1-\norm{R_2}_\infty}\,\norm{\widetilde M}_\infty .
\]
The same bound holds in any induced $p$-norm with $\norm{R_2}_p<1$:
$\norm{M-\widetilde M}_p\le\beta_p\coloneqq\frac{\norm{R_2}_p}{1-\norm{R_2}_p}\norm{\widetilde M}_p$.
Moreover, against the \emph{block-diagonal candidate}
$\Lambda=\operatorname{blockdiag}(P_1,\dots,P_q)$ of Section~\ref{sec:algo}, with the second
residual $R_1\coloneqq Y(AW-W\Lambda)$,
\begin{equation}\label{eq:blockdev}
  M-\Lambda=(I+R_2)^{-1}R_1,\qquad
  \norm{M-\Lambda}_2\;\le\;\beta_\Lambda\coloneqq\frac{\norm{R_1}_2}{1-\norm{R_2}_2}.
\end{equation}
The \emph{single} residual $\beta_\Lambda$ bounds the entire deviation of the collapsed
matrix from the block-diagonal candidate --- the inter-block leakage \emph{included} --- so
no separate coupling term is needed.
\end{lemma}

\begin{proof}
$W$ is nonsingular by Lemma~\ref{lem:neumann}, so $M=W^{-1}AW$ is a genuine
similarity and $\spec(A)=\spec(M)$. Using $W^{-1}=(I+R_2)^{-1}Y$,
\[
  M = W^{-1}AW = (I+R_2)^{-1}\,YAW = (I+R_2)^{-1}\widetilde M ,
\]
so $M-\widetilde M=\big((I+R_2)^{-1}-I\big)\widetilde M$ and Lemma~\ref{lem:neumann}
gives $\norm{M-\widetilde M}_p\le \frac{\norm{R_2}_p}{1-\norm{R_2}_p}\norm{\widetilde M}_p=\beta_p$.
For \eqref{eq:blockdev}, $M-\Lambda=W^{-1}(AW-W\Lambda)=(I+R_2)^{-1}Y(AW-W\Lambda)=(I+R_2)^{-1}R_1$,
and Lemma~\ref{lem:neumann} bounds $\norm{(I+R_2)^{-1}}_2\le 1/(1-\norm{R_2}_2)$.
\end{proof}

\begin{theorem}[Block enclosure]\label{thm:enclosure}
Suppose $\norm{R_2}_2<1$, and take the block-diagonal candidate
$\Lambda=\operatorname{blockdiag}(P_1,\dots,P_q)$, $P_i=\widetilde M_{\mathcal C_i\mathcal C_i}$,
with the single residual $\beta_\Lambda$ of \eqref{eq:blockdev}. Then
\[
  \spec(A)\ \subseteq\ \bigcup_{i=1}^{q}\Omega_i,
  \qquad \Omega_i\coloneqq\big\{\,z\in\C:\ \sigma_{\min}(P_i-zI)\le\beta_\Lambda\,\big\}.
\]
There is \emph{no} separate inter-block coupling term: with the block-diagonal $\Lambda$ as
the Miyajima candidate, the lone residual $\beta_\Lambda$ already bounds the full deviation
$\norm{M-\Lambda}_2$, the off-block leakage included. A singleton block
($\abs{\mathcal C_i}=1$) gives the disc $\Omega_i=\overline D(\widetilde M_{ii},\beta_\Lambda)$,
recovering a scalar Ger\v sgorin enclosure; the sharper per-row radius is
Remark~\ref{rem:sharpdisc}.
\end{theorem}

\begin{proof}
By Lemma~\ref{lem:deviation}, $\spec(A)=\spec(M)$ and $\norm{M-\Lambda}_2\le\beta_\Lambda$.
Since $\Lambda-zI$ is block diagonal,
$(\Lambda-zI)^H(\Lambda-zI)=\operatorname{blockdiag}\big((P_i-zI)^H(P_i-zI)\big)$, so its
singular values are the union of the blocks' and
$\sigma_{\min}(\Lambda-zI)=\min_i\sigma_{\min}(P_i-zI)$. Singular values are $1$-Lipschitz
(Weyl), so $\big|\sigma_{\min}(M-zI)-\sigma_{\min}(\Lambda-zI)\big|\le\norm{M-\Lambda}_2\le\beta_\Lambda$
for every $z$. If $\lambda\in\spec(M)$ then $\sigma_{\min}(M-\lambda I)=0$, whence
$\min_i\sigma_{\min}(P_i-\lambda I)=\sigma_{\min}(\Lambda-\lambda I)\le\beta_\Lambda$, i.e.
$\lambda\in\Omega_i$ for some $i$.
\end{proof}

\begin{remark}[Sharper discs for singletons]\label{rem:sharpdisc}
For a singleton block the uniform $\beta_\Lambda$ can be replaced by the per-row Ger\v sgorin
radius $\rho_i=\sum_{j\ne i}\abs{\widetilde M_{ij}}+\beta_\infty$ (scalar Ger\v sgorin on
$M$, recentred on $\widetilde M_{ii}$ with the row deviation absorbed into $\beta_\infty$): each
singleton disc shrinks from $\overline D(\widetilde M_{ii},\beta_\Lambda)$ to
$\overline D(\widetilde M_{ii},\rho_i)$, and the row/column form of
Section~\ref{sec:resolvent} sharpens it further. The block bound is the one that survives
when a cluster cannot be split; the per-row disc is its tight specialization where it can.
\end{remark}

\begin{proposition}[Block-Ger\v sgorin sharpening]\label{prop:blockgersg}
Theorem~\ref{thm:enclosure} charges every block the \emph{same} radius $\beta_\Lambda$,
the global deviation. It can be localized to each block's \emph{own} coupling. Write
$\widetilde N\coloneqq\widetilde M-\Lambda$ for the off-block part of $\widetilde M$ and, for
each block,
\begin{equation}\label{eq:blockrad}
  r_i\;\coloneqq\;\norm{\widetilde N_{\mathcal C_i,:}}_2+\beta_2,
  \qquad \beta_2=\tfrac{\norm{R_2}_2}{1-\norm{R_2}_2}\norm{\widetilde M}_2 ,
\end{equation}
with $\beta_2$ the $p=2$ slack of Lemma~\ref{lem:deviation}. Then
$\spec(A)\subseteq\bigcup_{i=1}^{q}\Omega_i'$ with
$\Omega_i'=\{z:\sigma_{\min}(P_i-zI)\le r_i\}$, and a connected component of overlapping
$\Omega_i'$ carries the summed multiplicities (Corollary~\ref{cor:clusters} applies verbatim).
\end{proposition}

\begin{proof}
Decompose the block row of the deviation as
$(M-\Lambda)_{\mathcal C_i,:}=(M-\widetilde M)_{\mathcal C_i,:}+\widetilde N_{\mathcal C_i,:}$.
Lemma~\ref{lem:deviation} gives $\norm{M-\widetilde M}_2\le\beta_2$, hence
$\norm{(M-\Lambda)_{\mathcal C_i,:}}_2\le r_i$. Now if $\lambda\in\spec(M)=\spec(A)$ with unit
eigenvector $x$, partition $x$ by the blocks and pick $i^\ast=\arg\max_i\norm{x_{\mathcal C_i}}_2$;
from $(M-\lambda I)x=0$, the $i^\ast$ block row reads
$(P_{i^\ast}-\lambda I)x_{\mathcal C_{i^\ast}}=-(M-\Lambda)_{\mathcal C_{i^\ast},:}\,x$, so
$\sigma_{\min}(P_{i^\ast}-\lambda I)\norm{x_{\mathcal C_{i^\ast}}}_2\le r_{i^\ast}\norm{x}_2
=r_{i^\ast}\norm{x_{\mathcal C_{i^\ast}}}_2$ (the last step since $x_{\mathcal C_{i^\ast}}$ has
maximal block norm and the bound is on the block \emph{row}; precisely,
$\norm{(M-\Lambda)_{\mathcal C_{i^\ast},:}\,x}_2\le\norm{(M-\Lambda)_{\mathcal C_{i^\ast},:}}_2\norm{x}_2$).
Dividing, $\sigma_{\min}(P_{i^\ast}-\lambda I)\le r_{i^\ast}$. The multiplicity count is the
homotopy argument of Corollary~\ref{cor:clusters} with the radii $r_i$ fixed.
\end{proof}

This is the block analogue of Feingold--Varga: each disc is inflated only by its own
off-block-row coupling $\norm{\widetilde N_{\mathcal C_i,:}}_2$ (plus the small inverse-slack
$\beta_2$), never by the worst block's. Where the coupling is uneven --- many weakly-coupled
singletons beside one strongly-coupled near-defective block --- the singleton discs collapse
to nearly $\overline D(\widetilde M_{ii},\beta_2)$ while only the bad block keeps a large
radius. In the worst block $r_i\le\norm{\widetilde N}_2+\beta_2$, recovering the global bound;
it is never looser. The sharpening matters for the disc \emph{radii}; it changes the certified
\emph{count} (Corollary~\ref{cor:clusters}) only when the coupling, not the eigenvalue gap, is
the binding constraint --- e.g.\ the GKW operator of Section~\ref{sec:gkw} has off-block
coupling $\norm{N}_2\approx2\times10^{-23}$, already orders below its gaps, so its localized
and global counts coincide ($89$ free singletons either way).

\begin{remark}[The candidate eigenvalues drop out]\label{rem:Dfree}
The enclosure of Theorem~\ref{thm:enclosure} depends only on the basis $W$ (through
$\widetilde M=YAW$, $R_2=YW-I$): the candidate diagonal $D$ does not appear. The block
centres are the diagonal blocks $P_i$ of $\widetilde M$ --- the (block) Rayleigh-quotient
refined eigenvalue estimates --- and $\beta_\Lambda$ is independent of $D$. Refinement
(Section~\ref{sec:algo}) drives the residual $R_1=Y(AW-W\Lambda)$, hence $\beta_\Lambda$,
toward the $\kappa u$ floor, sharpening the enclosure.
\end{remark}

\begin{corollary}[Cluster count]\label{cor:clusters}
Under the hypothesis of Theorem~\ref{thm:enclosure}, if a subcollection of blocks
$\{\mathcal C_i\}_{i\in S}$ has $\bigcup_{i\in S}\Omega_i$ connected and disjoint from the
remaining regions $\Omega_{i'}$, then it contains exactly $\sum_{i\in S}\abs{\mathcal C_i}$
eigenvalues of $A$, counted with algebraic multiplicity.
\end{corollary}

\begin{proof}
Apply Theorem~\ref{thm:enclosure} along the homotopy $M_t=\Lambda+t(M-\Lambda)$,
$t\in[0,1]$: each $M_t$ has $\norm{M_t-\Lambda}_2\le\beta_\Lambda$, so $\spec(M_t)\subseteq\bigcup_i\Omega_i$
throughout, with the $\Omega_i$ fixed. At $t=0$ the spectrum is $\bigcup_i\spec(P_i)$,
placing $\abs{\mathcal C_i}$ eigenvalues in $\Omega_i$; the eigenvalues vary continuously in
$t$ without leaving the regions, so a component disjoint from the rest retains its total
count $\sum_{i\in S}\abs{\mathcal C_i}$.
\end{proof}

This is the certified output: the regions $\Omega_i$ of Theorem~\ref{thm:enclosure},
grouped by Corollary~\ref{cor:clusters} into connected components each carrying a known
eigenvalue count. A near-defective or strongly coupled cluster, which the geometry of
Section~\ref{sec:geom} forbids the refinement from separating, is exactly one block $P_i$:
its region $\Omega_i=\{\sigma_{\min}(P_i-zI)\le\beta_\Lambda\}$ is a single connected set carrying
its full multiplicity $\abs{\mathcal C_i}$ --- the enclosure never tries to resolve the
ill-posed interior.

\section{From the enclosure to a resolvent bound, and experiments}\label{sec:resolvent}

The certificate of Section~\ref{sec:rigor} is a statement about the \emph{spectrum};
the same data bound the \emph{resolvent} $(B-zI)^{-1}$ cheaply, away from a small
exclusion set. Since $B-zI = W(M-zI)W^{-1}$ with $M=W^{-1}BW$,
\begin{equation}\label{eq:transport}
  \sigma_{\min}(B-zI)\;\ge\;\frac{\sigma_{\min}(M-zI)}{\kappa_2(W)},
  \qquad \kappa_2(W)=\norm{W}_2\norm{W^{-1}}_2 .
\end{equation}
The collapsed $M=W^{-1}BW$ deviates from the block-diagonal candidate
$\Lambda=\operatorname{blockdiag}(P_1,\dots,P_q)$ by $\norm{M-\Lambda}_2\le\beta_\Lambda$,
the \emph{single} Miyajima residual \eqref{eq:blockdev}. Block diagonality gives
$\sigma_{\min}(\Lambda-zI)=\min_j\sigma_{\min}(P_j-zI)$, and $\sigma_{\min}(\cdot)$ is
$1$-Lipschitz (Weyl), so $\sigma_{\min}(M-zI)\ge\min_j\sigma_{\min}(P_j-zI)-\beta_\Lambda$.
With \eqref{eq:transport} this is the \emph{primary} bound of the paper:

\begin{lemma}[Block resolvent floor]\label{lem:blockres}
With $\beta_\Lambda$ the single residual of \eqref{eq:blockdev} and $\kappa_2(W)$ from $R_2$,
\[
  \sigma_{\min}(B-zI)\;\ge\;
  \frac{\max\big(0,\ \min_j\sigma_{\min}(P_j-zI)-\beta_\Lambda\big)}{\kappa_2(W)} .
\]
Equivalently $\norm{(B-zI)^{-1}}_2\le\kappa_2(W)\big/\big(\min_j\sigma_{\min}(P_j-zI)-\beta_\Lambda\big)$
where the denominator is positive --- i.e.\ outside the enclosure $\bigcup_i\Omega_i$. There
is no separate inter-block coupling term: $\beta_\Lambda$ subsumes it. The residual
$\beta_\Lambda$ and $\kappa_2(W)$ are computed once from the certificate; per $z$ the cost is one small
$\sigma_{\min}(P_j-zI)$ per block (a closed form for $\abs{\mathcal C_j}\le2$, a verified
SVD of $P_j-zI$ otherwise), and a singleton block contributes only the far field
$\abs{z-\widetilde M_{jj}}$. This is where the architecture pays off: the non-normality of a
defective cluster lives in $\sigma_{\min}(P_j-zI)$, which tracks $\abs{z-\lambda}^m$
honestly, \emph{not} in $\kappa_2(W)$ --- the latter carries only the benign inter-block
conditioning ($\approx1/\sin\theta_{\mathrm{inter}}$).
\end{lemma}

Lemma~\ref{lem:blockres} subtracts the \emph{global} $\beta_\Lambda$ from every block. The
same block-Ger\v sgorin localization that sharpens the enclosure (Proposition~\ref{prop:blockgersg})
sharpens the resolvent: charge each block only its own coupling radius $r_j$ of
\eqref{eq:blockrad}.

\begin{proposition}[Localized block-dominance floor]\label{prop:blockdom}
Let $s_j(z)=\sigma_{\min}(P_j-zI)$ and $r_j$ as in \eqref{eq:blockrad}. If
$\tau(z)^2\coloneqq\sum_j\big(r_j/s_j(z)\big)^2<1$ then
\[
  \sigma_{\min}(M-zI)\;\ge\;\Big(\min_j s_j(z)\Big)\big(1-\tau(z)\big),
  \qquad
  \sigma_{\min}(B-zI)\;\ge\;\frac{(\min_j s_j(z))(1-\tau(z))}{\kappa_2(W)} .
\]
\end{proposition}

\begin{proof}
Factor $M-zI=(\Lambda-zI)\big(I+(\Lambda-zI)^{-1}(M-\Lambda)\big)$, valid since each
$s_j(z)>0$ makes $\Lambda-zI$ invertible. The matrix $G=(\Lambda-zI)^{-1}(M-\Lambda)$ has
block rows $G_{\mathcal C_j,:}=(P_j-zI)^{-1}(M-\Lambda)_{\mathcal C_j,:}$, so
$\norm{G_{\mathcal C_j,:}}_2\le r_j/s_j(z)$ (Proposition~\ref{prop:blockgersg}). For any
block-row-partitioned $G$, $\norm{G}_2^2=\norm{\sum_j G_{\mathcal C_j,:}^H G_{\mathcal C_j,:}}_2
\le\sum_j\norm{G_{\mathcal C_j,:}}_2^2\le\tau(z)^2$. Hence
$\sigma_{\min}(I+G)\ge1-\norm{G}_2\ge1-\tau(z)$ and, with
$\sigma_{\min}(\Lambda-zI)=\min_j s_j(z)$ and $\sigma_{\min}(XY)\ge\sigma_{\min}(X)\sigma_{\min}(Y)$,
the first bound follows; \eqref{eq:transport} gives the second.
\end{proof}

The engine returns the maximum of Lemma~\ref{lem:blockres} and
Proposition~\ref{prop:blockdom} --- both rigorous lower bounds. The dominance floor improves
on the global one wherever the coupling is uneven across blocks (it replaces $\beta_\Lambda$
by the small $r_j$ of weakly-coupled blocks); the global floor wins near a block's own
spectrum, where some $s_j(z)\to0$ pushes $\tau(z)\ge1$ and the factorization degenerates.
On well-separated spectra ($\beta_\Lambda$ and all $r_j$ far below the gaps) the two agree to
rounding, as on the gallery of Table~\ref{tab:gallery}; the gain is reserved for the
coupling-bound regime.

\paragraph{Block-diagonal estimate vs.\ the exact block SVD.} Both Lemma~\ref{lem:blockres}
and Proposition~\ref{prop:blockdom} read $\sigma_{\min}$ off the \emph{diagonal} blocks $P_j$
and divide by the transport $\kappa_2(W)$ of \eqref{eq:transport}: they certify
$\sigma_{\min}(B-zI)$ through the block-diagonal surrogate, at the cost of one small
$\sigma_{\min}(P_j-zI)$ per block. When $B$ is small enough to factor a verified SVD directly,
$\sigma_{\min}(B-zI)$ is instead certified \emph{exactly} by a single $\texttt{svdbox}(B-zI)$
per $z$, with \emph{no} $\kappa_2$ factor --- the SVD sees the true non-normality of the whole
block rather than the block-diagonal proxy plus its inter-block conditioning. For a captured
Krylov block of moderate size $m=\dim B$ (hundreds) this costs $O(m^3)$ per sample
(milliseconds) and recovers exactly the $\kappa_2$ that the block-diagonal floor concedes. On
the GKW operator with an $m=32$ Krylov block, $\kappa_2(W)\approx4.3$, and replacing the
block-diagonal estimate by the exact $\texttt{svdbox}$ lifts the certified $\sigma_{\min}$ on
the inter-eigenvalue circle from $0.024$ to the true $0.075$ --- turning a marginal gap circle
into a comfortable one. The per-block layer is therefore reserved for blocks too large for one
SVD, where its $O(\sum_j m_j^3)\ll O(m^3)$ cost is worth the $\kappa_2$ price; for the moderate
blocks of the Krylov gap certificate the exact SVD is both cheap and tight, and is the default.

\paragraph{The cheap all-singletons layer.} For a separable spectrum the blocks are
singletons, $\sigma_{\min}(P_j-zI)=\abs{z-\widetilde M_{jj}}$, and the block floor is a
scalar gap that the \emph{recentred Ger\v sgorin gaps} of $\widetilde M=YBW$ sharpen in both
norms (read up to $\norm{M-\widetilde M}_p\le\beta_p$,
$\beta_p=\tfrac{\norm{R_2}_p}{1-\norm{R_2}_p}\norm{\widetilde M}_p$) --- the per-row
($\ell^\infty$) and per-column ($\ell^1$) gap:
\[
  g_\infty(z)=\min_i\Big(\abs{z-\widetilde M_{ii}}-\!\sum_{j\ne i}\abs{\widetilde M_{ij}}-\beta_\infty\Big),
  \qquad
  g_1(z)=\min_j\Big(\abs{z-\widetilde M_{jj}}-\!\sum_{i\ne j}\abs{\widetilde M_{ij}}-\beta_1\Big).
\]
Where these are positive, diagonal dominance gives $\norm{(M-zI)^{-1}}_\infty\le 1/g_\infty(z)$
and $\norm{(M-zI)^{-1}}_1\le 1/g_1(z)$, and Riesz--Thorin interpolates them to $\ell^2$:
$\sigma_{\min}(M-zI)\ge\sqrt{g_1(z)\,g_\infty(z)}$ --- the $O(n)$-per-$z$, SVD-free
specialization of Lemma~\ref{lem:blockres}. With \eqref{eq:transport}:

\begin{lemma}[Free resolvent floor, row/column form]\label{lem:freeres}
\[
  \sigma_{\min}(B-zI)\;\ge\;
  \frac{\sqrt{\max(0,g_1(z))\cdot\max(0,g_\infty(z))}}{\kappa_2(W)} .
\]
The centres $\widetilde M_{ii}$, the row/column off-diagonal sums and the slacks
$\beta_\infty,\beta_1$ are precomputed once ($O(n^2)$, \emph{no} SVD), so the bound costs
$O(n)$ per $z$. The gaps are $z$-LOCAL --- the binding row and column change with $z$ ---
so the bound adapts to where $z$ sits instead of paying one global far-field everywhere;
$\kappa_2(W)$ is the only global factor. For a normal $B$ ($\kappa_2=1$, $\widetilde M$
diagonal) it returns the exact $\operatorname{dist}(z,\sigma(B))$.
\end{lemma}

\begin{proposition}[Where the floor is undefined]\label{prop:grey}
The floor is positive exactly outside the union of recentred discs
\[
  \mathcal G=\bigcup_i \overline D(\widetilde M_{ii},\rho_i^{\mathrm{row}})\,\cup\,
                       \overline D(\widetilde M_{ii},\rho_i^{\mathrm{col}}),\qquad
  \rho_i^{\mathrm{row}}=\!\sum_{j\ne i}\abs{\widetilde M_{ij}}+\beta_\infty,\quad
  \rho_i^{\mathrm{col}}=\!\sum_{j\ne i}\abs{\widetilde M_{ji}}+\beta_1,
\]
a $z$-LOCAL exclusion set --- the row/column off-diagonal mass about each $\widetilde
M_{ii}$. A coupled / non-normal cluster shows up directly as a wide disc (a large
row/column tail); a separable spectrum has tiny tails and $\mathcal G$ shrinks to the
eigenvalues. ($g_\infty$ or $g_1\le0$ exactly inside $\mathcal G$.)
\end{proposition}

\subsection*{Numerical experiments}

We run the full VBD pipeline (Section~\ref{sec:algo}) --- Schur, Newton--Rump separation,
block-diagonal candidate, within-block orthonormalization, Miyajima certificate, block
resolvent floor (Lemma~\ref{lem:blockres}) --- on the standard gallery: a normal matrix, the
Grcar matrix, a defective Jordan block, two tight non-normal clusters, and a weakly graded
chain. Table~\ref{tab:gallery} records the block structure, the certificate quantities, and
the floor's accuracy against the true $\sigma_{\min}(A-zI)$ on a $25\times25$ box grid.

\begin{table}[h]\centering
\begin{tabular}{lrlrrrr}
\toprule
matrix & $n$ & blocks & $\kappa_2$ & $\norm{R_2}_2$ & $\beta_\Lambda$ & far / blkSVD \\
\midrule
normal       & 16 & $16{\times}1$ & $1.0$ & $1.1\mathrm{e}{-}14$ & $1.3\mathrm{e}{-}13$ & $100\%$ / $1.00$ \\
Grcar        & 32 & $1{\times}32$ & $1.0$ & $2.2\mathrm{e}{-}14$ & $5.7\mathrm{e}{-}13$ & $4\%$ / $1.00$ \\
Jordan       &  8 & $1{\times}8$  & $1.0$ & $4.4\mathrm{e}{-}16$ & $2.0\mathrm{e}{-}15$ & $40\%$ / $1.00$ \\
two clusters &  6 & $3,2,1$       & $3.1$ & $5.2\mathrm{e}{-}15$ & $6.8\mathrm{e}{-}14$ & $86\%$ / $0.37$ \\
chain        & 12 & $12{\times}1$ & $3.9$ & $1.3\mathrm{e}{-}15$ & $9.2\mathrm{e}{-}14$ & $100\%$ / $0.36$ \\
\bottomrule
\end{tabular}
\caption{The VBD pipeline on the gallery. \emph{blocks}: the merge-only clustering keeps a
globally non-normal (Grcar) or defective (Jordan) spectrum as a single block, while a
separable spectrum splits into singletons. $\kappa_2$ is \emph{benign everywhere}
($\le3.9$): the non-normality is confined to the blocks, not the basis. \emph{far / blkSVD}:
fraction of the grid certified by the cheap far-field floor $\abs{z-c_j}-\norm{N_j}$, and
the median tightness (floor$/$true $\sigma_{\min}$) of the per-$z$ block-SVD floor --- which
is \emph{exact} ($1.00$) on the non-normal Grcar and Jordan blocks. At every one of the
$625$ grid points the certified floor lies below the true $\sigma_{\min}$ (rigor holds).}
\label{tab:gallery}
\end{table}

\begin{figure}[h]\centering
  \begin{minipage}{0.46\textwidth}\centering
    \includegraphics[width=\linewidth]{bench/figures/free_resolvent_jordan8.png}\\
    {\small (a) defective Jordan: grey disc of radius $=$ row-tail $1$}
  \end{minipage}\hfill
  \begin{minipage}{0.46\textwidth}\centering
    \includegraphics[width=\linewidth]{bench/figures/free_resolvent_grcar32.png}\\
    {\small (b) Grcar, eigenvector frame: no grey, but loose by $\kappa_2\approx10^5$}
  \end{minipage}
  \caption{The \emph{cheap free-floor layer}, illustrating its two limitations --- both
  removed by the VBD block bound of Table~\ref{tab:gallery}. (a) On the defective Jordan
  block the floor leaves a grey disc of radius $=$ row-tail $1$ (Proposition~\ref{prop:grey};
  cyan dashed circle on the grey boundary), which the VBD per-$z$ block SVD fills
  \emph{exactly}. (b) On Grcar the \emph{eigenvector}-frame floor is loose everywhere by
  $\kappa_2\approx10^5$, the failure the VBD frame removes by keeping Grcar as one block with
  $\kappa_2=1$. Heatmap $=\log_{10}\norm{(B-zI)^{-1}}_2$ upper bound, grey $=\mathcal G$
  (undefined), red $\times$ eigenvalues.}
  \label{fig:gallery}
\end{figure}

The table shows the architecture doing its job: the non-normality never reaches
$\kappa_2$. Two cases.
\begin{itemize}
  \item \emph{Globally non-normal or defective} (Grcar, Jordan). The merge-only clustering
    cannot separate the spectrum, so it keeps the whole matrix as one block and the basis is
    the \emph{orthonormal Schur frame}, $\kappa_2=1$. The cheap far-field floor is then loose
    --- it sees the block as one fat disc of radius $\norm{N_j}$, certifying only $4\%$
    (Grcar) / $40\%$ (Jordan) of the grid --- but the per-$z$ block SVD of $P_j-zI$ recovers
    $\sigma_{\min}$ \emph{exactly} (tightness $1.00$), tracking the defective
    $\abs{z-\lambda}^m$ decay. The non-normality lives in the block, where the SVD handles
    it, not in the prefactor.
  \item \emph{Separable} (normal, chain, the loosely-coupled clusters). The spectrum splits
    into singletons or small blocks, $\kappa_2$ stays benign ($1.0$--$3.9$), and the cheap
    far-field floor already certifies most of the grid ($86$--$100\%$) at the exact gap.
\end{itemize}

\paragraph{Why not the eigenvector frame?} The same gallery on the \emph{eigenvector} basis
is a cautionary tale. Diagonalizing Grcar --- globally non-normal but with a well-separated
spectrum --- forces an exponentially ill-conditioned basis, $\kappa_2\approx9.8\times10^{4}$,
and the transported floor $\sqrt{g_1g_\infty}/\kappa_2$ is loose by that factor
\emph{everywhere} (at $z=3+3i$, true $\sigma_{\min}=1.2$ versus $2.4\times10^{-5}$): with the
spectrum separated, no block absorbs the non-normality, so it has nowhere to go but the
prefactor. The VBD frame removes exactly this failure mode --- by \emph{not} diagonalizing a
block it cannot cleanly separate, it keeps $\kappa_2=1$ and pushes the difficulty into
$\sigma_{\min}(P_j-zI)$, where one $m\times m$ SVD resolves it. The price is the per-$z$
block SVD, but only inside the kept blocks' small discs, and the blocks are exactly as large
as the local non-normality demands (here, $32$ for Grcar, $8$ for Jordan, $1$ everywhere
else).

\subsection*{Application: certifying the Gauss--Kuzmin--Wirsing spectrum}\label{sec:gkw}

The free floor is built for exactly the regime where a \emph{coarse} resolvent bound
suffices. In the certified spectral approximation of transfer operators
(Nisoli, \texttt{arXiv:2602.19435}), a finite section $A_K$ of an operator $L$ is tied to
$L$ by a \emph{small-gain} (resolvent-bridge) condition: on a circle $\Gamma_j$ encircling
only the candidate eigenvalue $\lambda_j$,
\[
  \alpha_j \;=\; \varepsilon_K\,\sup_{z\in\Gamma_j}\norm{(A_K-zI)^{-1}}_2 \;<\; 1
  \quad\Longrightarrow\quad L \text{ has exactly one eigenvalue inside } \Gamma_j,
\]
$\varepsilon_K$ bounding the truncation error. The costly ingredient is
$\sup_{\Gamma_j}\norm{(A_K-zI)^{-1}}_2$ --- classically a verified SVD of $A_K-zI$ at every
contour sample. Lemma~\ref{lem:freeres} supplies it from a \emph{single} enclosure of
$A_K$ at $O(\#\text{clusters})$ per $z$. Because $\varepsilon_K$ is tiny, $\alpha_j<1$
carries dozens of orders of headroom, so even the $\kappa_2$-inflated free bound closes
the bridge --- the loose prefactor is irrelevant here.

We take the $K=256$ Gauss--Kuzmin--Wirsing matrix ($n=257$, $\varepsilon_{256}\approx
5.6\times10^{-45}$). Its spectrum decays geometrically to $0$ ($\lambda_k\sim(2/3)^k$), so
the near-$0$ eigenvalues are too close to separate at any finite precision: the merge-only
clustering correctly keeps them as a \emph{single block} and peels the well-separated
leading eigenvalues off as singletons. Table~\ref{tab:gkw} runs the pipeline at increasing
precision, comparing the VBD frame against the eigenvector frame.

\begin{table}[h]\centering
\begin{tabular}{lrrrrr}
\toprule
frame / precision & blocks & near-$0$ blk & $\kappa_2$ & $\norm{R_2}_\infty$ & free \\
\midrule
eigenvector, Float64 & $257{\times}1$ & ---   & $\sim10^{14}$ (artifact)   & ---                  & $2$  \\
eigenvector, 256-bit & $257{\times}1$ & ---   & $2.8\times10^{6}$          & $\sim10^{-70}$       & $89$ \\
eigenvector, 512-bit & $257{\times}1$ & ---   & $2.0\times10^{11}$         & $\sim10^{-150}$      & $77$ \\
VBD, 256-bit         & $90$           & $168$ & $1.3\times10^{2}$          & $1.5\times10^{-74}$  & $54$ \\
VBD, 512-bit         & $158$          & $100$ & $1.4\times10^{2}$          & $1.2\times10^{-151}$ & $49$ \\
\bottomrule
\end{tabular}
\caption{The GKW $K=256$ matrix ($n=257$) at increasing precision. The \emph{eigenvector}
frame forces every eigenvalue into a $1\times1$ block, so it must diagonalize the
near-defective near-$0$ tail: $\kappa_2$ is huge and \emph{grows} with precision
($2.8\times10^6\!\to\!2.0\times10^{11}$) and the free count is \emph{non-monotone}
($89\!\to\!77$) --- the signature of a fragile certificate. The \emph{VBD} frame keeps the
unsplittable tail as one block ($168$ at 256-bit, shrinking to $100$ at 512-bit as more
leading eigenvalues resolve and peel off), so $\kappa_2$ stays \emph{benign and stable}
($\approx1.4\times10^2$) and $\norm{R_2}$ falls to the precision floor. \emph{free}:
eigenvalues whose small-gain $\alpha_j<1$ closes.}
\label{tab:gkw}
\end{table}

Two honest readings. For the small-gain bridge --- whose headroom is enormous
($\varepsilon_{256}\sim10^{-45}$) --- even the $\kappa_2$-inflated eigenvector floor closes
$89$ contours for free at 256-bit (the loose prefactor is irrelevant against $40$ orders of
headroom); this is the cheapest route and the original purpose of the free floor. But that
count rests on a $\kappa_2$ that \emph{blows up} with precision --- the eigenvector frame
straining to separate eigenvalues that are not well-posedly separable --- so it is not
robust: at 512-bit it \emph{drops} to $77$. The VBD frame trades count for robustness:
$\kappa_2$ benign and flat, the near-$0$ tail certified as a \emph{single block} (the
well-posed object --- its interior is the ill-posed problem one must not try to resolve),
and $54$ genuinely-isolated singletons closing the bridge. Finer resolution erodes the
isolation count ($54\!\to\!49$, more resolved eigenvalues crowd the gaps) even as the
\emph{decomposition} sharpens monotonically ($\norm{R_2}\!\to\!10^{-151}$, block
$168\!\to\!100$, $\kappa_2$ flat). Either way the cost is decisive: one $O(n^3)$ enclosure
$+$ a contour scan at $O(\#\text{blocks})$ per $z$, against a per-$z$ verified SVD of
$A_K-zI$ at every contour sample --- $\sim\!400\times$ cheaper in Float64, the gap widening
sharply in extended precision.

\begin{figure}[h]\centering
  \includegraphics[width=\textwidth]{bench/figures/gkw_free_resolvent_zoom0.png}
  \caption{The \emph{eigenvector}-frame free floor near $0$ for the $K=256$ GKW matrix
  (Table~\ref{tab:gkw}, eigenvector rows), Float64 (left) vs BigFloat (right); grey $=$
  undefined ($z\in\mathcal G$), heatmap $=\log_{10}$ of the certified resolvent upper bound,
  $\times$ eigenvalues, circles the isolating contours $\Gamma_j$ (green $=$ small-gain
  closed for free, red $=$ escalate). Float64 leaves a $\sim\!10^{-2}$ grey hole about the
  small eigenvalues ($2$ certified); the BigFloat candidate resolves them and erases the
  hole, certifying $89$ for free --- at the cost of the precision-unstable $\kappa_2$ of
  Table~\ref{tab:gkw}, which the VBD frame avoids.}
  \label{fig:gkw}
\end{figure}

\section{Matrix-free operators: certifying the far block by matvec}\label{sec:matfree}

Up to here $A$ is dense and $W$ is formed and stored. When $A$ is available only through
the maps $x\mapsto Ax$, $x\mapsto A^{*}x$ --- a sparse matrix, or an implicitly defined
operator such as a discretized transfer operator or a differential operator --- the dense
decomposition of Section~\ref{sec:algo} can still be run on the small subspace of
eigenvalues near a target $z$, but not on the whole spectrum. This section certifies
$\sigma_{\min}(A-zI)$ in that regime by splitting the work into a dense \emph{captured}
block and a matrix-free \emph{far} block, the latter handled by a projected Lanczos
process whose basis conditioning is certified \emph{a posteriori from scalar coefficients},
with no stored basis. The single residual $\norm{R_2}_2<1$ that drives the dense theory
reappears here and, remarkably, certifies at once that the Lanczos frame is invertible
\emph{and} that it is complete.

\subsection*{The two-block split}

Run the VBD of Section~\ref{sec:algo} on the cluster of eigenvalues near $z$, producing a
block-orthonormal frame $W_{\mathcal C}\in\C^{n\times m}$ ($m$ small) with
$W_{\mathcal C}^{*}W_{\mathcal C}=I_m$ exactly. Because the columns are orthonormal,
\[
  P_{11}=W_{\mathcal C}W_{\mathcal C}^{*},\qquad P_{22}=I-P_{11}
\]
are exact orthogonal projectors --- no approximate inverse is needed, block-orthonormality
makes $Y_{\mathcal C}=W_{\mathcal C}^{*}$ exact --- and $A$ has the exact $2\times2$ block
form in the orthogonal splitting $\C^{n}=\operatorname{ran}P_{11}\oplus\operatorname{ran}P_{22}$,
\[
  A=\begin{pmatrix}G_{11}&G_{12}\\G_{21}&G_{22}\end{pmatrix},\quad
  G_{11}=W_{\mathcal C}^{*}AW_{\mathcal C},\ \ G_{12}=P_{11}AP_{22},\ \
  G_{21}=P_{22}AP_{11},\ \ G_{22}=P_{22}AP_{22}.
\]
The far block acts on $\operatorname{ran}P_{22}$ by $G_{22}x=P_{22}AP_{22}x$ --- two
matvecs and two rank-$m$ updates --- and no entry of $G_{22}$ is ever formed. Fix $z$ and
put $M_z:=(G_{22}-zI)^{*}(G_{22}-zI)\succeq0$ on $\operatorname{ran}P_{22}$. The
compression is clean: the shift $-zI$ is block-diagonal and $P_{22}^2=P_{22}$, so
\[
  M_z=P_{22}(A-zI)^{*}P_{22}(A-zI)P_{22}=(G_{22}-zI)^{*}(G_{22}-zI)
\]
on $\operatorname{ran}P_{22}$ --- the coupling $G_{12},G_{21}$ does not contaminate $M_z$,
so $\sigma_{\min}(G_{22}-zI)=\sqrt{\lambda_{\min}(M_z)}$ exactly, and the goal is a
certified lower bound on $\lambda_{\min}(M_z)$ from matvecs on $M_z$ alone (two enclosed
matvecs on $A,A^{*}$ per application).

\subsection*{The Paige recurrence: basis conditioning without a stored basis}

Lanczos on the Hermitian $M_z$ builds an orthonormal Krylov basis of
$\operatorname{ran}P_{22}$ by a three-term recurrence; in finite precision the computed
$\bar q_j$ satisfy
\begin{equation}\label{eq:lanczos-fp}
  \bar\beta_j\bar q_{j+1}=M_z\bar q_j-\bar\alpha_j\bar q_j-\bar\beta_{j-1}\bar q_{j-1}+f_j,
\end{equation}
with $f_j$ the local defect, and lose orthogonality: the Gram entries
$\omega_{ij}=\bar q_i^{*}\bar q_j$ drift from $\delta_{ij}$, and the basis-conditioning
residual is $R_2=\bar Q_k^{*}\bar Q_k-I$ with entries $\omega_{ij}$ ($i\ne j$). A direct
evaluation of $R_2$ would need all $k$ stored $n$-vectors. Paige's recurrence removes that
need: left-multiplying \eqref{eq:lanczos-fp} by $\bar q_i^{*}$, using $M_z=M_z^{*}$ and
\eqref{eq:lanczos-fp} at index $i$, gives for $1\le i<j$
\begin{equation}\label{eq:paige}
  \bar\beta_j\,\omega_{i,j+1}
  =\bar\beta_i\,\omega_{i+1,j}+(\bar\alpha_i-\bar\alpha_j)\,\omega_{ij}
   +\bar\beta_{i-1}\,\omega_{i-1,j}-\bar\beta_{j-1}\,\omega_{i,j-1}+\xi_{ij},
  \quad \xi_{ij}=\bar q_i^{*}f_j-f_i^{*}\bar q_j,
\end{equation}
propagating one anti-diagonal of the Gram matrix from the previous two using only the
scalar coefficients and the defects. Evaluated in interval arithmetic --- absolute values
in the numerator, the lower endpoint $\underline{\bar\beta_j}$ in the denominator, and the
Cauchy--Schwarz bound
$\abs{\xi_{ij}}\le\overline{\norm{\bar q_i}}\,\overline{\norm{f_j}}+\overline{\norm{f_i}}\,\overline{\norm{\bar q_j}}$
for the residual term --- it yields an inductive entrywise bound $W_{ij}\ge\abs{\omega_{ij}}$,
seeded by $\omega_{ij}=\delta_{ij}$ on the first two anti-diagonals. Every input is a scalar:
$\bar\alpha_j,\bar\beta_j$, the defect norms $\overline{\norm{f_j}}$ (a byproduct of each
enclosed Lanczos step, whose two matvecs are already required), and the certified vector
norms $\overline{\norm{\bar q_j}}$. No $n$-vector is stored. Finally, since $\abs{R_2}\le W$
entrywise and the $2$-norm is monotone under entrywise magnitude domination,
\begin{equation}\label{eq:r2bound}
  \norm{R_2}_2\;\le\;\norm{W}_2\;\le\;\sqrt{\norm{W}_1\,\norm{W}_\infty},
\end{equation}
a rigorous bound from the streamed scalars. This is genuinely a posteriori --- no error
model, no a-priori constant. Classical a-priori loss-of-orthogonality bounds propagate
worst-case rounding through the same recurrence and routinely overshoot by orders of
magnitude; \eqref{eq:paige} runs on the \emph{realized} $\bar\alpha_j,\bar\beta_j,\overline{\norm{f_j}}$,
so it reflects the actual conditioning of the run and can fall below $1$ where the a-priori
envelope does not. (It takes absolute values in the propagation, so it does not exploit
sign cancellation; the gain is the realized coefficients, not cancellation.)

\subsection*{The full basis: one residual certifies invertibility and completeness}

Run Lanczos to a \emph{full} basis, $k=\dim\operatorname{ran}P_{22}$. The single test
$\norm{R_2}_2<1$ then does double duty.

\begin{lemma}[Full-basis certificate]\label{lem:fullbasis}
Let $\bar Q_k$ have $k=\dim\operatorname{ran}P_{22}$ columns in $\operatorname{ran}P_{22}$,
and suppose \eqref{eq:paige}--\eqref{eq:r2bound} prove $\norm{R_2}_2\le\rho<1$,
$R_2=\bar Q_k^{*}\bar Q_k-I$. Then $\bar Q_k^{*}\bar Q_k=I+R_2\succ0$, so $\bar Q_k$ has
full column rank and its columns form a basis of $\operatorname{ran}P_{22}$. Hence
$\bar Q_k$ is a nonsingular frame on $\operatorname{ran}P_{22}$ with $Y=\bar Q_k^{*}$,
$R_2=Y\bar Q_k-I$, and Lemma~\ref{lem:deviation} applies with $M=M_z$, $W=\bar Q_k$,
$\Lambda=T_k$ (the computed tridiagonal): every eigenvalue of $M_z$ lies within
$\beta^{(22)}=\norm{R_1}_2/(1-\rho)$ of $\spec(T_k)$, where $R_1=\bar Q_k^{*}(M_z\bar Q_k-\bar Q_kT_k)$.
\end{lemma}

Completeness is what removes the failure mode of a Krylov \emph{compression}
($k\ll\dim$): there, an eigenvector orthogonal to the $k$-dimensional Krylov space is
invisible to the matvec data, and a small eigenvalue of $M_z$ --- a near-$z$ eigenvalue of
$G_{22}$ --- can hide in the unseen complement, producing a floor that exceeds the true
value. With a full basis and $\rho<1$ there is no unseen complement: full column rank forces
the columns to span $\operatorname{ran}P_{22}$, so the enclosure of Lemma~\ref{lem:deviation}
runs over the \emph{whole} spectrum of $M_z$. The same scalar test certifies both that the
frame is invertible and that nothing escaped it --- the Neumann condition $\norm{R_2}<1$ is
simultaneously a nonsingularity and a completeness certificate.

\subsection*{The matvec floor}

\begin{proposition}[Matvec floor on the far block]\label{prop:matfloor}
Under Lemma~\ref{lem:fullbasis}, let $\theta_{\min}=\lambda_{\min}(T_k)$ be the smallest
Ritz value (a certified small symmetric eigensolve, $O(k)$), bounded below by
$\underline\theta_{\min}$, and let the streamed eigen-residual of Lemma~\ref{lem:fullbasis}
give $\norm{R_1}_2\le\rho_1$. Then
\[
  \lambda_{\min}(M_z)\;\ge\;\underline\theta_{\min}-\beta^{(22)},\qquad
  \beta^{(22)}=\frac{\rho_1}{1-\rho},
\]
so $\sigma_{\min}(G_{22}-zI)\ge\sqrt{\max(0,\ \underline\theta_{\min}-\beta^{(22)})}$, a
certified floor from matvecs on $M_z$ alone --- no stored basis, no entries of $G_{22}$, no
factorization.
\end{proposition}

\begin{proof}
By Lemma~\ref{lem:fullbasis}, $\spec(M_z)\subseteq\bigcup_i\overline D(\theta_i,\beta^{(22)})$.
Since $M_z$ is Hermitian, $\theta_{\min}\ge\lambda_{\min}(M_z)$, and the eigenvalue of $M_z$
nearest $\theta_{\min}$ is within $\beta^{(22)}$, giving
$\lambda_{\min}(M_z)\ge\theta_{\min}-\beta^{(22)}$; replace $\theta_{\min}$ by its lower
endpoint and use $\sigma_{\min}(G_{22}-zI)=\sqrt{\lambda_{\min}(M_z)}$.
\end{proof}

As in the dense theory the conditioning $\kappa_2$ does not enter the eigenvalue floor
(similarity preserves the spectrum); were the resolvent \emph{norm} itself wanted, it is
certified from the streamed data alone through $\norm{\bar Q_k}_2^2=\norm{I+R_2}_2\le1+\rho$,
giving $\overline\kappa\le(1+\rho)/(1-\rho)$, benign for $\rho<1$. The eigen-residual $R_1$
streams from the same scalars: from $M_z\bar Q_k=\bar Q_kT_k+\bar\beta_k\bar q_{k+1}e_k^{*}+F_k$,
the Ritz residual of $(\theta_i,y_i)$, $y_i=\bar Q_k\bar s_i$, is
$r_i=\bar\beta_k(\bar s_i)_k\bar q_{k+1}+F_k\bar s_i$; at a full basis $\bar\beta_k\to0$ and
$r_i$ reduces to the accumulated local defects $F_k\bar s_i$, so $\norm{R_1}_2$ is bounded
from the streamed $\overline{\norm{f_j}}$ and the small $T_k$ eigendata, e.g.\ by
$\sqrt{\norm{R_1}_1\norm{R_1}_\infty}$ accumulated in one further deterministic pass (a
running-max float and one $n$-vector, $O(n)$ storage).

\begin{remark}[Soundness in the far-from-spectrum regime]\label{rem:zfar}
The certificate is one-directional, and the splitting is designed so that it is used only
where it fires. Orthogonality is lost in Lanczos when a Ritz value \emph{converges}: when
$z$ is far from $\spec(G_{22})$, $\lambda_{\min}(M_z)=\sigma_{\min}(G_{22}-zI)^2$ is lifted
away from $0$ with no bottom cluster, the bottom Ritz value does not converge early, the
divisor $\underline{\bar\beta_j}$ in \eqref{eq:paige} stays healthy, and $\norm{R_2}_2<1$
holds out to the full basis --- the sound, firing regime. A near-$z$ eigenvalue drives
$\lambda_{\min}(M_z)\to0$, collapses a $\bar\beta_j$, and pushes $\norm{R_2}_2\ge1$, so the
certificate declines. By construction the near-$z$ eigenvalues are captured in $G_{11}$, so
the far block is always evaluated at $z$ far from $\spec(G_{22})$: a decline there signals
insufficient working precision (raise it, or enlarge the captured block), never a missed
near-$z$ eigenvalue --- unlike the general case, the failure here carries no spectral
ambiguity. A \emph{proved} bound is never wrong: it certifies the far block with no
eigenvalue overlooked.
\end{remark}

\subsection*{Assembly: the coupling and the combined floor}

The off-diagonal blocks $G_{12},G_{21}$ are the inter-block coupling; they are computed
directly, not absorbed into a slack. Since $P_{11}=W_{\mathcal C}W_{\mathcal C}^{*}$,
\[
  \norm{G_{21}}_2=\norm{P_{22}AW_{\mathcal C}}_2
  =\norm{AW_{\mathcal C}-W_{\mathcal C}(W_{\mathcal C}^{*}AW_{\mathcal C})}_2,
\]
the block Rayleigh--Ritz (right-invariance) residual of $W_{\mathcal C}$, formed from the
$m$ matvecs $AW_{\mathcal C}$ already needed by the VBD; likewise
$\norm{G_{12}}_2=\norm{A^{*}W_{\mathcal C}-W_{\mathcal C}(W_{\mathcal C}^{*}A^{*}W_{\mathcal C})}_2$
from $A^{*}W_{\mathcal C}$. For non-normal $A$ the two differ and \emph{both} are needed.
Both are $z$-independent (the shift is block-diagonal), so the coupling
\[
  \gamma:=\max\!\bigl(\norm{G_{12}}_2,\norm{G_{21}}_2\bigr)\le\overline\gamma
\]
is certified \emph{once}, before the $z$-sweep, through the $m\times m$ Gram matrices
$C^{*}C$ (cheap geometric-mean bound, or a verified $m\times m$ SVD).

The basis $[\,W_{\mathcal C},W_\perp\,]$ ($W_\perp$ any orthonormal basis of
$\operatorname{ran}P_{22}$) is orthonormal, so the similarity to the $2\times2$ block form
is unitary and preserves $\sigma_{\min}$ --- no $\kappa_2$ enters the outer split. With
$E=\bigl(\begin{smallmatrix}0&G_{12}\\G_{21}&0\end{smallmatrix}\bigr)$, $E^{*}E$ is
block-diagonal, so $\norm{E}_2=\max(\norm{G_{12}}_2,\norm{G_{21}}_2)=\gamma$ exactly.
Weyl's inequality on $A-zI=D+E$, $D=\diag(G_{11}-zI,\,G_{22}-zI)$, gives the combined floor.

\begin{theorem}[Matrix-free pseudospectrum floor]\label{thm:matfree}
If the dense VBD certifies $G_{11}$ (Lemma~\ref{lem:blockres}) and the Paige certificate of
Proposition~\ref{prop:matfloor} proves $\norm{R_2}_2<1$ for the full far-block basis, then
\[
  \sigma_{\min}(A-zI)\;\ge\;
  \min\!\bigl(\sigma_{\min}(G_{11}-zI),\,\sigma_{\min}(G_{22}-zI)\bigr)-\overline\gamma,
\]
with $\sigma_{\min}(G_{11}-zI)\ge\min_{j}\sigma_{\min}(P_j-zI)-\beta_\Lambda^{(11)}$ the
dense VBD floor (the intra-$G_{11}$ non-normality lives inside this term, not in the outer
split) and $\sigma_{\min}(G_{22}-zI)$ the matvec floor of
Proposition~\ref{prop:matfloor}. The right-hand side is assembled from one certified block
SVD per small block of $G_{11}$, $O(k)$ projected matvecs with $O(k^2)$ scalar a-posteriori
work for $G_{22}$, and the single $z$-independent coupling $\overline\gamma$ --- no global
matrix is ever formed, and no a-priori assumption on the spectrum is made.
\end{theorem}

\section{Application: certified range of a polynomial}\label{sec:polyrange}

The spectral enclosure of Theorem~\ref{thm:enclosure} is, at bottom, a verified root
finder: the eigenvalues of a companion matrix are the roots of a polynomial, and the discs
enclose \emph{all} of them. This turns the engine into a rigorous range enclosure --- given
$p(x)=\sum_{k=0}^d c_k x^k$ with real $c_k$ and an interval $[a,b]$, certify an interval
containing $\{p(x):x\in[a,b]\}$ --- which makes a clean illustration of where the
eigenvalue route is the right tool and where it is not. We contrast two routes, both
rigorous on the same midpoint--radius backend.

\paragraph{Route 1: branch-and-bound.} Subdivide $[a,b]$, evaluate $p$ in ball arithmetic
on each subinterval (Horner, outward-rounded), and keep a worklist; split the box whose
enclosure drives the current extreme, and prune boxes that provably cannot hold the
minimum (resp.\ maximum) --- the interval Moore--Skelboe iteration. It needs no derivative,
degrades gracefully, but its cost is set by the requested tolerance: the box overestimation
is $O(h)$ in the box width (linear), so reaching a value-bracket of width $\tau$ takes
$O(1/\tau)$ boxes, and on an \emph{oscillatory} $p$ --- many extrema of nearly equal value
--- pruning fails and the cost explodes. Certifying the range of the Chebyshev polynomial
$T_{10}$ on $[-1,1]$ to width $10^{-3}$ already takes $\approx 9.6\times10^{4}$ ball
evaluations ($0.3$\,s); to width $10^{-6}$, $\approx 3\times10^{6}$ ($\approx 4$\,min).

\paragraph{Route 2: critical points via the companion enclosure.} For $p\in C^1$ the
extrema of $p$ on $[a,b]$ are attained at the endpoints or at interior stationary points
$p'(x)=0$. Form $p'$, certify the roots of its companion matrix with the engine
(Theorem~\ref{thm:enclosure}), keep every cluster disc that meets the real axis inside
$[a,b]$, and hull the ball-evaluations of $p$ at the endpoints and over those stationary
enclosures.

\begin{proposition}[Soundness of the critical-point route]\label{prop:polyrange}
Let $\mathcal D=\bigcup_j D(c_j,r_j)$ be the certified discs for the roots of $p'$
(Theorem~\ref{thm:enclosure}: every root of $p'$ lies in $\mathcal D$). For each $j$ with
$\abs{\operatorname{Im}c_j}\le r_j$ let $I_j=[\operatorname{Re}c_j-s_j,\operatorname{Re}c_j+s_j]\cap[a,b]$,
$s_j=\sqrt{r_j^2-(\operatorname{Im}c_j)^2}$, be the disc's real trace clipped to $[a,b]$
(empty when the disc clears the axis). Then with $V$ the ball-hull of
$\{p(a),p(b)\}\cup\{p(I_j)\}$ (each evaluated in ball arithmetic),
$\{p(x):x\in[a,b]\}\subseteq V$.
\end{proposition}
\begin{proof}
A real stationary point $x_\star\in[a,b]$ is a real root of $p'$, hence an eigenvalue of the
companion, hence in some $D(c_j,r_j)$; being real it lies in that disc's real trace, so
$x_\star\in I_j$ and $p(x_\star)\in p(I_j)\subseteq V$. The endpoints are in $V$ by
construction. Since the global extrema of $p$ over $[a,b]$ are attained among
$\{a,b\}\cup\{x_\star\}$ and ball evaluation is inclusion-isotonic, both extremal values lie
in $V$; $V$ is an interval, so it contains the whole range. No stationary point is missed
because the discs enclose \emph{all} roots of $p'$ --- the same completeness that
Corollary~\ref{cor:clusters} guarantees for the spectrum.
\end{proof}

The decisive feature is that route 2's accuracy is \emph{decoupled from any tolerance}: one
eigenproblem returns the range to the width of the root enclosure, no matter how tight. On
$T_{10}$ above it returns the exact $[-1,1]$ in $0.5$\,ms, against route 1's $0.3$\,s--$4$\,min.

\paragraph{How bad are the timings for route 2?} They are not the binding constraint --- the
\emph{conditioning} is. A single call is one degree-$(d{-}1)$ companion eigendecomposition
plus one Miyajima certification, i.e.\ $O(d^3)$ ball products, independent of the target
accuracy; the float-Newton refinement is negligible and helps only the clustered case, not
the well-separated one. In \texttt{Float64} this is sub-millisecond through degree $\approx16$
and a few milliseconds at degree $32$ (Table~\ref{tab:polyrange}, top). But the companion is
the textbook non-normal matrix, and on the Chebyshev family $\kappa_2$ grows by roughly an
order of magnitude every two degrees; around degree $32$--$64$ the \texttt{Float64}
certificate's Neumann residual $\norm{R_2}_\infty$ exceeds $1$ and the engine \emph{soundly
refuses} (it cannot certify the basis nonsingular). Recovering the range past that wall needs
extended precision, and \emph{that} is where the time goes: \texttt{BigFloat} ball products
carry a large constant, so at $256$ bits degree $32$ costs $\approx0.75$\,s and degree $64$
$\approx6$\,s; worse, the precision itself must scale with the degree --- $256$ bits is
already exhausted by degree $96$, where the enclosure blows up (width $\sim10^{19}$) --- so
once precision-bound the cost is effectively super-cubic. The honest summary: route 2 is
extremely fast where it is well-conditioned (low-to-moderate degree, well-separated extrema),
and where it is ill-conditioned the limit is not the wall-clock of a single solve but the
precision that solve must run at.

\begin{table}[h]\centering
\begin{tabular}{rrrr}
\toprule
degree $d$ & $\kappa_2$(companion) & max disc radius & time \\
\midrule
\multicolumn{4}{l}{\emph{\texttt{Float64}}}\\
$4$  & $4.9$               & $3.5\mathrm{e}{-}15$ & $0.06$\,ms \\
$8$  & $1.4\times10^{2}$   & $4.7\mathrm{e}{-}14$ & $0.12$\,ms \\
$16$ & $1.4\times10^{5}$   & $3.4\mathrm{e}{-}11$ & $0.36$\,ms \\
$32$ & $1.5\times10^{11}$  & $2.7\mathrm{e}{-}5$  & $1.5$\,ms \\
$64$ & \multicolumn{3}{l}{Neumann fail ($\norm{R_2}_\infty\approx44\ge1$): refuses, soundly} \\
\midrule
\multicolumn{4}{l}{\emph{\texttt{BigFloat}, $256$ bits}}\\
$16$ & $1.4\times10^{5}$   & $1.3\mathrm{e}{-}70$ & $82$\,ms \\
$32$ & $1.5\times10^{11}$  & $2.6\mathrm{e}{-}60$ & $0.75$\,s \\
$64$ & $2.3\times10^{23}$  & $9.9\mathrm{e}{-}38$ & $6.4$\,s \\
$96$ & $3.9\times10^{35}$  & \multicolumn{2}{l}{width $\sim10^{19}$: $256$ bits exhausted ($30$\,s)} \\
\bottomrule
\end{tabular}
\caption{Route 2 on the Chebyshev polynomials $T_d$ over $[-1,1]$ (true range $[-1,1]$,
$d{-}1$ interior extrema --- the adversarial conditioning case for the companion). Per-call
cost is $O(d^3)$ and tolerance-independent, so the timings are small until the companion's
conditioning forces extended precision. In \texttt{Float64} the certificate is lost between
degree $32$ and $64$; \texttt{BigFloat} restores it, at a cost dominated by the working
precision, which must itself grow with the degree.}
\label{tab:polyrange}
\end{table}

\section{Remarks}\label{sec:remarks}

\begin{remark}[Why no inverse, and conditioning in one residual]
The only place $W^{-1}$ could enter is through $M=W^{-1}AW$; Lemma~\ref{lem:deviation}
replaces it by the computed $\widetilde M=YAW$ at the cost of the single slack $\beta$,
which is controlled by $\norm{R_2}_\infty$ alone. The conditioning of the eigenbasis
never appears in the enclosure (it is scale-invariant); it is relevant only as a
reported diagnostic $\kappa_2(W)=\norm{W}_2\norm{W^{-1}}_2$, for which a tight bound is
obtained from a verified singular value decomposition rather than from the cheaper
$\sqrt{\norm{\cdot}_1\norm{\cdot}_\infty}$ estimate, which overestimates by $O(\sqrt n)$
on orthonormal bases.
\end{remark}

\begin{remark}[The asymmetry that licenses a numerical refinement]
Refinement and clustering are decided in floating point and never certified. The
single certification at the end is self-correcting: an over-optimistic split produces
discs that still overlap and is silently re-merged by Corollary~\ref{cor:clusters},
while an attempt to separate genuinely coincident eigenvalues drives the Newton step
\eqref{eq:correction} to blow up, ill-conditions $W$, and is caught as
$\norm{R_2}_\infty\ge1$. Either way the output is never wrong.
\end{remark}

\begin{remark}[Full rigor in ball arithmetic]
To make Theorem~\ref{thm:enclosure} a theorem about \emph{floating-point} computation,
the products $YW$ and $YAW$ are evaluated in midpoint--radius (ball) arithmetic, so
$\widetilde M$ and $R_2$ are ball matrices that rigorously contain the exact products;
the row sums $\sum_{j\ne i}\abs{\widetilde M_{ij}}$ and the norm $\norm{R_2}_\infty$ are
then certified upper bounds. The proof above is unchanged: $\beta$ and the $\rho_i$
become rigorous over-estimates, and the discs still enclose $\spec(A)$.
\end{remark}

\section*{To do}
\begin{itemize}
  \item Localized refinement: confine the per-cluster work to the $m\times m$
    sub-block and the $n\times m$ subspace columns ($O(nm^2)$ rather than $O(n^3)$).
  \item Quantitative refinement bound: relate the converged radii to the eigenvector
    conditioning and the working precision (the $u^{1/m}$ floor for an $m$-fold
    cluster).
  \item Extended-precision local re-solve for the first-order--obstructed clusters of
    Remark~\ref{rem:exp} (the bases that stall for \emph{both} corrections): redo the
    $m\times m$ block in higher precision (\S4d, \texttt{BigFloat}) rather than seeking
    a different first-order step.
  \item Numerical examples and the comparison with Rump's \texttt{verifyeigall}.
  \item Matrix-free validation (Section~\ref{sec:matfree}): a ball Lanczos on the far block
    $G_{22}$ with the streamed Paige bound \eqref{eq:paige}, measuring how the certified
    $\norm{R_2}_2$ grows toward $1$ with the basis dimension and with $\operatorname{dist}(z,\spec G_{22})$,
    on the Gauss--Kuzmin--Wirsing operator.
\end{itemize}

\end{document}
